% BHU PYQ -- Manually transcribed & error-corrected
% Paper: GLB-101: Elementary Physical and Structural Geology
% Exam: B.Sc. (Hons.) Semester I Examination, 2016-17

\documentclass[11pt,a4paper]{article}
\usepackage[a4paper,top=2.5cm,bottom=2.5cm,left=2.8cm,right=2.8cm,headheight=14pt]{geometry}
\usepackage{amsmath,amssymb}
\usepackage[T1]{fontenc}
\usepackage[utf8]{inputenc}
\usepackage{lmodern,microtype,fancyhdr,enumitem,hyperref,lastpage,parskip}

\hypersetup{colorlinks=true,urlcolor=black,linkcolor=black,
  pdftitle={GLB-101 Elementary Physical and Structural Geology -- BHU B.Sc. Sem I 2016-17}}
\pagestyle{fancy}\fancyhf{}
\renewcommand{\headrulewidth}{0.4pt}\renewcommand{\footrulewidth}{0.4pt}
\fancyhead[L]{\small   \textit{Banaras Hindu University}}
\fancyhead[R]{\small   \textit{GLB-101: Elementary Physical and Structural Geology}}
\fancyfoot[C]{\small Page  \thepage\ of \pageref{LastPage}}


\newlist{parts}{enumerate}{2}
\setlist[parts,1]{label=(\alph*),leftmargin=*,itemsep=5pt,topsep=3pt}

\newlist{romanparts}{enumerate}{2}
\setlist[romanparts,1]{label=(\roman*),leftmargin=*,itemsep=5pt,topsep=3pt}

\newcommand{\pts}[1]{\hfill   \textit{[#1]}}


\begin{document}

\begin{center}
  {\large\bfseries BANARAS HINDU UNIVERSITY}\\[2pt]
  {
\normalsize B.Sc. (Hons.) --- Semester I Examination, 2016-17}\\[6pt]
  \rule{0.8\linewidth}{0.5pt}\\[6pt]
  {\large\bfseries Paper No. GLB-101: Elementary Physical and Structural Geology}\\[4pt]
  {
\normalsize Time: 3 Hours \hfill Maximum Marks: 70}\\[2pt]
  {\small Ref. No.: BS/Sem I/166}
\end{center}
\rule{\linewidth}{0.4pt}
\smallskip



\noindent   \textit{Instructions: Attempt any   \textbf{five} questions, selecting at least   \textbf{two} questions each from Section A and Section B. Section C is compulsory. All questions carry equal marks (14 marks each).}
\bigskip

\bigskip


\noindent {\large\bfseries SECTION --- A}
\rule{\linewidth}{0.4pt}\smallskip




\noindent   \textbf{Question 1.} Discuss in detail with diagrams the internal constitution of the Earth. \pts{14}

\medskip


\noindent   \textbf{Question 2.} What do you understand by earthquakes? Discuss their causes, classification, and geographical distribution. \pts{14}

\medskip


\noindent   \textbf{Question 3.} Define fold. What are the different parts of a fold? With the help of diagrams, give the classification of folds based on: \pts{14}

\begin{parts}
  \item arching direction
  \item relative ages of beds
  \item plunge of the hinge lines
\end{parts}

\medskip


\noindent   \textbf{Question 4.} Write explanatory notes on the following: \pts{$2  \times 7 = 14$}

\begin{parts}
  \item Unconformity
  \item Forms of igneous rocks
\end{parts}

\bigskip


\noindent {\large\bfseries SECTION --- B}
\rule{\linewidth}{0.4pt}\smallskip




\noindent   \textbf{Question 5.} Discuss the following: \pts{$2  \times 7 = 14$}

\begin{parts}
  \item Convection currents in the Earth's core and production of the geomagnetic field
  \item Geological time scale
\end{parts}

\medskip


\noindent   \textbf{Question 6.} Write the salient differences between the following: \pts{$4  \times 3.5 = 14$}

\begin{parts}
  \item Normal fault and reverse fault
  \item Nappe and window
  \item Outlier and inlier
  \item Overlap and offlap
\end{parts}

\medskip


\noindent   \textbf{Question 7.} Write short notes on   \textbf{any four} of the following: \pts{$4  \times 3.5 = 14$}

\begin{parts}
  \item Dip and strike
  \item Clinometer compass
  \item Relative and absolute methods of age determination
  \item Age of the Earth
  \item Briefly discuss the origin of the atmosphere
  \item Concordant and discordant bodies
\end{parts}

\bigskip


\noindent {\large\bfseries SECTION --- C (Compulsory)}
\rule{\linewidth}{0.4pt}\smallskip




\noindent   \textbf{Question 8.} Write the answers to the following in two or three sentences, or fill in the blanks as the case may be: \pts{14}

\begin{parts}
  \item What is anticlinorium? \pts{1}
  \item What is sill? \pts{1}
  \item What is master joint? \pts{1}
  \item What is thrust fault? \pts{1}
  \item What is fan fold? \pts{1}
  \item What type of dip does a dome structure have? \pts{1}
  \item What is outcrop? \pts{1}
  \item Earthquake magnitude is measured on the \underline{\hspace{3cm}} scale and intensity is measured on the \underline{\hspace{3cm}} scale. \pts{2}
  \item The most explosive volcanic eruption is of the \underline{\hspace{3cm}} type and the least explosive is of the \underline{\hspace{3cm}} type. \pts{2}
  \item The Big Bang happened \underline{\hspace{3cm}} years ago and the Milky Way galaxy was formed \underline{\hspace{3cm}} years ago. \pts{2}
  \item The volcanic belt around the Pacific Ocean is known as the \underline{\hspace{4cm}}. \pts{1}
\end{parts}

\vfill\rule{\linewidth}{0.4pt}

\begin{center}\small
\href{https://bhu-exam.vercel.app/}{Click here to visit our website}\\[4pt]
\href{https://me-aryan.vercel.app/}{Click here to visit our contributor}\\[4pt]
\href{https://quashan-me.vercel.app/}{Click here to visit our contributor}
\end{center}
\end{document}
